%%% Research-paper draft: outcome-anchored segmentation of coding-agent trajectories.
%%% Template: the author's acmart setup (TrajAnalysis.tex / reprot.tex).
%%% For a double-blind submission, replace the class options with
%%%   \documentclass[sigconf,review,anonymous]{acmart}
%%% and remove the \settopmatter / \setcopyright lines.
\documentclass[sigconf,nonacm,10pt]{acmart}

\usepackage{microtype}
\usepackage{graphicx}
\usepackage{booktabs}
\usepackage{multirow}
\usepackage{amsmath}
\usepackage{balance}
\usepackage{xcolor}
\usepackage{algorithm}
\usepackage{algpseudocode}

\settopmatter{printacmref=false, printccs=false, printfolios=true}
\setcopyright{none}
\renewcommand\footnotetextcopyrightpermission[1]{}
\pagestyle{plain}

%%% Draft helpers. Every \TODO marks a number or claim that does not exist yet.
\newcommand{\TODO}[1]{\textcolor{red}{[TODO: #1]}}
\newcommand{\tool}{\textsc{TrajSeg}}   % working name; rename freely
\newcommand{\code}[1]{\texttt{#1}}

\begin{document}

\title{Where Did the Repair Go Wrong? Outcome-Anchored Segmentation of
       Coding-Agent Trajectories}

\author{Mustahid Hasan}
\affiliation{%
  \institution{William \& Mary}
  \city{Williamsburg}
  \state{Virginia}
  \country{USA}}
\email{mhasan02@wm.edu}

\renewcommand{\shortauthors}{Hasan}

\begin{abstract}
Coding agents leave long trajectories of tool calls, and researchers read them to
find out why a repair failed. Existing trajectory viewers either show the raw step
list or tag each step with a phase from fixed rules over tool and command names.
None groups steps into the contiguous, multi-step units an analyst actually reasons
about, such as ``steps 4--9 were one failed attempt to reproduce the bug'', and none
validates its labels against human judgement. We present \tool{}, which normalises
trajectories from any agent into the Agent Trajectory Interchange Format (ATIF),
describes each tool call by what it does rather than by its tool name, and segments
the trajectory by \emph{outcome-anchored optimal partitioning}: an exact dynamic
program over an objective whose evidence is the result of each run (did the test
pass, fail, or never run?) and the trajectory's state (before or after the first
source edit). The same failing test is thus a successful reproduction before the fix
and a failed fix after it. Segments are typed by a decision table shared with a
human-rater codebook, grouped into seven repair purposes by Viterbi decoding, and
flagged for failure and inefficiency in six coordinated views. We evaluate \tool{} on
150 trajectories of Claude Code, Codex and Gemini CLI on 50 Android program-repair
bugs, against a blind two-rater reference set and re-implemented baselines from
existing tools. \TODO{headline RQ1--RQ2 numbers once the reference set is scored.}
\end{abstract}

\keywords{LLM agents, agent trajectories, automated program repair, segmentation,
visualization, event abstraction}

\maketitle

%% ===================================================================== 1
\section{Introduction}
\label{sec:intro}

LLM-based coding agents such as Claude Code, Codex and Gemini CLI repair bugs by
interleaving reasoning with tool calls: they read and search code, edit files, run
builds and tests, and react to the output~\cite{yao2023react,yang2024sweagent}. The
record of this process, the \emph{trajectory}, is the main evidence for
understanding why an agent succeeded or failed. Trajectories are also long and
heterogeneous. In our corpus a single repair takes up to 117 tool calls, a single
tool result can exceed 24 million characters of build output, and each agent logs in
its own format.

Researchers therefore build tools to read trajectories. We surveyed 19 of them
(Section~\ref{sec:background}). They show trajectories as timelines, Gantt charts or
collapsible step lists, and several claim ``multiple levels of abstraction''. Their
higher levels are of two kinds: \emph{structural} groupings that follow the log
(spans, API calls, sub-agents), or a \emph{phase tag} on each step assigned by fixed
rules over tool or command names, as in Graphectory~\cite{graphectory2026} and Agent
ATO~\cite{agentato2026}. Neither matches how an analyst explains a run. An analyst
says ``actions 1--3 tried to run the failing test but the project would not compile;
actions 6--7 fixed the build configuration and the test then failed as expected''.
That sentence names \emph{contiguous multi-step units}, each with a \emph{goal} and
an \emph{outcome}. The outcome depends on context: the same failing test is a success
before the fix (the bug is reproduced) and a failure after it (the fix does not
work). No existing tool produces such units, and none of the phase taggers reports
agreement with human labels.

This paper presents \tool{}, which recovers these units automatically and makes the
failing and wasteful ones visible. Its key design decisions are:
\begin{itemize}
  \item \textbf{Agent-agnostic input.} Trajectories are converted to
        ATIF~\cite{atif2026} and each tool call is described as an \emph{action
        frame}: operation, object and effect. Frames are inferred from the call's
        arguments and the program a shell command runs, not from tool names, so a
        new agent needs no hand-written table.
  \item \textbf{Outcome-anchored segmentation as an optimisation problem.}
        Segmentation maximises an explicit objective whose evidence is each run's
        result and the trajectory's pre-/post-fix state. An exact dynamic program
        finds the optimum, following optimal partitioning and semi-Markov
        segmentation~\cite{jackson2005optimal,sarawagi2004semicrf}. The objective's
        terms map one-to-one to the rules of a human-rater codebook.
  \item \textbf{A shared specification for tool and raters.} Segment types come
        from a decision table that both the tool and the human raters apply, so
        disagreement measures either tool error or rubric ambiguity, never two
        different definitions.
\end{itemize}

\paragraph{Contributions.}
\begin{itemize}
  \item \textbf{C1, outcome-anchored segmentation}
        (Sections~\ref{sec:seg}--\ref{sec:purpose}): a formulation of trajectory
        segmentation as constrained optimal partitioning with result and state
        evidence, and Viterbi decoding of repair purposes.
  \item \textbf{C2, agent-agnostic action frames} (Section~\ref{sec:frames}):
        classification of tool calls from argument signatures and executed
        programs, tested by leave-one-agent-out.
  \item \textbf{C3, surfacing failure and inefficiency} (Section~\ref{sec:surface}):
        segment-level flags and loops, shown in six coordinated views.
  \item \textbf{C4, evaluation} (Sections~\ref{sec:design}--\ref{sec:results}):
        150 ATIF trajectories with harness outcomes, a two-rater codebook and
        reference set, and a comparison against re-implemented baselines from
        existing tools. All are released.
\end{itemize}

%% ===================================================================== 2
\section{Background and Motivation}
\label{sec:background}

\subsection{Agent trajectories and ATIF}
A trajectory is the sequence of an agent's inference steps. Each step may hold
reasoning text, one or more tool calls, and their results. Agents log trajectories
in incompatible formats: Claude Code as SDK message objects, Codex as a rollout
event stream, Gemini CLI as a JSON stream interleaved with stderr. ATIF~\cite{atif2026}
is a JSON schema for such trajectories, maintained by the Harbor project, with
explicit steps, tool calls, observations and embedded sub-agent trajectories. We use
ATIF v1.7 as the single input format. Our converters are checked by two independent
validators (ours and Harbor's) and by a retention audit that searches the converted
file for every text-bearing block of the raw log.

\subsection{Motivating example}
\label{sec:example}
Figure~\ref{fig:timeline} shows Gemini CLI on bug 18 of our corpus, a fragment-manager
crash in an Android app. The harness reports \emph{fail}. A step list shows 37 tool
calls and nothing more. Read as segments, the run tells a story:
\begin{itemize}
  \item actions 1--3 try to run the bug-revealing test and fail on a compile error;
  \item actions 6--7 edit \code{build.gradle}, after which the test runs and fails,
        so the bug is reproduced;
  \item a 14-action localization sweep follows;
  \item action 22 edits the fragment, and the test passes, as it does twice more
        (actions 29 and 32);
  \item the agent then edits \code{build.gradle} again, breaks the build, edits it
        once more, and stops without running anything.
\end{itemize}
The agent's final message claims success. The failure is explained by the last
segment: an unverified build-configuration change after a failing build. A per-step
phase tag labels actions 33--37 ``environment'' or ``patching'' and loses both the
outcome and the fact that nothing was verified.

\begin{figure*}[t]
  \centering
  \includegraphics[width=\textwidth]{figs/traj-timeline.png}
  \caption{Timeline view of Gemini CLI on bug 18 (harness: fail), from the current
  prototype. Rows: purposes; segments (fill = outcome); actions (fill = result); the
  failure rail (shape = failure kind). The \textsc{claim}$\times$ badge marks a run
  whose final message claims success while the harness reports failure.}
  \Description{A horizontal timeline with purpose bands, segment blocks coloured by
  outcome, per-action result markers and a failure rail for 37 actions.}
  \label{fig:timeline}
\end{figure*}

\subsection{What existing tools offer}
Table~\ref{tab:tools} summarises the ten tools of our 19-tool survey that are closest
to this work. Three gaps recur. No tool partitions a run into contiguous multi-step
units. Every semantic label is either manual (TraceView~\cite{traceview2026}) or a
per-step tag from fixed rules (Graphectory~\cite{graphectory2026}, Agent
ATO~\cite{agentato2026}, Mindwalk~\cite{mindwalk}). Only Inspect
Scout~\cite{inspectscout} validates automatic labels against humans, and it labels
whole transcripts, not their structure. ProcGrep~\cite{oderinwale2026procgrep} induces
procedures by byte-pair encoding over 16 hand-defined atoms. On our corpus this is a
co-occurrence model with no notion of outcome, and we use it as a baseline.

\begin{table*}[t]
\centering
\small
\caption{Closest existing trajectory tools (from our 19-tool survey). ``Segments'':
contiguous multi-step units with boundaries. ``Validated'': automatic labels
compared with human labels.}
\label{tab:tools}
\resizebox{\textwidth}{!}{%
\begin{tabular}{@{}llllll@{}}
\toprule
\textbf{Tool} & \textbf{Levels in a run} & \textbf{Semantic labels} & \textbf{Segments} & \textbf{Validated} & \textbf{Outcome analysis} \\
\midrule
Agent Trajectory Explorer~\cite{desmond2025explorer} & TAO turns / raw & none & no & no & no \\
TraceView~\cite{traceview2026} & step / TAR & manual & no & no & outcome tag \\
SeaView~\cite{seaview2025} & experiment / instance / run & none & no & no & experiment diff \\
Graphectory~\cite{graphectory2026} & action + phase tag & rule per command type & no & no & GSP mining, Sankey \\
Agent ATO~\cite{agentato2026} & interaction + tag & rule tags (4 kinds) & no & no & excluded \\
Inspect Scout~\cite{inspectscout} & span timeline & LLM scanners & no & yes (P/R/F1) & filter by success \\
Mindwalk~\cite{mindwalk} & 160 buckets & read vs.\ edit & no & no & no \\
AgentDiff~\cite{agentdiff} & flat DAG & none & no & no & baseline gating \\
ProcGrep~\cite{oderinwale2026procgrep} & atoms $\rightarrow$ BPE procedures & frequency & tokens only & no & log-odds pass/fail \\
AgentDiagnose~\cite{ou2025agentdiagnose} & whole run & LLM scores & no & $r=0.57$ ($n=30$) & data filtering \\
\midrule
\textbf{\tool{} (ours)} & action $\rightarrow$ segment $\rightarrow$ purpose & result- and state-anchored & \textbf{yes} & \textbf{yes} (RQ1) & per segment and run \\
\bottomrule
\end{tabular}}
\end{table*}

%% ===================================================================== 3
\section{Approach}
\label{sec:approach}

\tool{} turns a raw log into a three-level hierarchy:
\begin{itemize}
  \item \textbf{actions} (one tool call each);
  \item \textbf{segments} (contiguous runs of actions with one immediate goal, a
        type and an outcome);
  \item \textbf{purposes} (contiguous runs of segments forming one stage of repair).
\end{itemize}
Every action belongs to exactly one segment and every segment to exactly one purpose.

\subsection{Action frames}
\label{sec:frames}
Each tool call becomes a frame $(o, b, r)$:
\begin{itemize}
  \item $o$, the \emph{operation}: read, search, list, write, execute, wait,
        delegate, plan, fetch or communicate;
  \item $b$, the \emph{object}: application source, test, toolchain, harness,
        environment, version control or external;
  \item $r$, the \emph{effect}: the run's result (Section~\ref{sec:verdict}).
\end{itemize}

\paragraph{Operation.} The operation is resolved in a fixed order, and the step that
decided it is recorded:
\begin{enumerate}
  \item \textbf{Format-level facts.} A call with a sub-agent reference is a delegation;
        a call carrying a session id returned by an earlier call is a wait.
  \item \textbf{Argument signature.} Old and new text or a patch body indicate a write;
        a pattern or query, a search; only a file path, a read; a command string, an
        execution.
  \item \textbf{The executed program.} A shell command is split into simple commands,
        respecting quotes, \code{\$(\ldots)} substitutions and heredoc bodies.
        Each part is classified by the program it runs, looking through wrappers such
        as \code{env} and \code{sudo} and recursing into \code{bash -c} scripts. A
        compound command takes the highest-priority category of its parts; for
        example, \code{cd app \&\& ./gradlew test | tail} runs tests.
  \item \textbf{Tool names}, only to break ties.
  \item \textbf{An LLM}, for a tool none of the above resolves: one call per
        \emph{distinct tool}, with a cached answer that a person checks.
\end{enumerate}

\paragraph{Object.} The object is read from the files involved:
\begin{itemize}
  \item test files follow each ecosystem's conventions (\code{test/},
        \code{androidTest/}, \code{*Test.java}, \code{test\_*.py});
  \item toolchain files are build and dependency configuration;
  \item harness files are named in the task prompt.
\end{itemize}

\paragraph{Compatibility.} The 15 mid-level categories used by earlier work and by our
viewer (e.g.\ \code{run\_test}, \code{edit\_file}) are a fixed function of the frame.
In our corpus, 52\% of tool calls are shell commands, already classified by the
program they run. The remaining 48\% use 18 agent-specific tool names, which the
argument signatures replace.

\subsection{Result parsing}
\label{sec:verdict}
The exit code does not tell whether a test failed: the corpus harness exits 0 when the
UI test fails, so an exit-code-based failure flag misses test failures for two of the
three agents. \tool{} therefore reads each execution's output and assigns a result
(\emph{pass}, \emph{fail} or \emph{none}) and, for failures, a kind:
\begin{itemize}
  \item \emph{test}: the test executed and failed;
  \item \emph{build}: the project's code does not compile;
  \item \emph{env}: the toolchain or sandbox stopped the run;
  \item \emph{exit}: any other non-zero exit.
\end{itemize}
Generic parsers for JUnit, Gradle and pytest output are separated from a per-benchmark
adapter, as SWE-bench separates per-repository log parsers~\cite{jimenez2024swebench}.
A pass requires test evidence. ``BUILD SUCCESSFUL'' or exit code 0 after a test command
gives \emph{none}, because a truncating pipe such as \code{| head} can hide the result.
A \emph{decisive run} is a test run with a pass or fail result, or a failed build.

\subsection{Outcome-anchored segmentation}
\label{sec:seg}
Let $a_1,\dots,a_n$ be the actions and $A$ the \emph{anchor}, the first edit of
application source. A segmentation $S$ partitions $1..n$ into contiguous spans; a
\emph{cut} $c$ separates $a_c$ from $a_{c+1}$. We choose
\begin{equation}
S^{*} = \arg\max_{S}\; \sum_{s\in S}\varphi(s) \;+\; \sum_{c\in\mathrm{cuts}(S)} b(c).
\label{eq:objective}
\end{equation}

\paragraph{Span score.}
\begin{equation}
\varphi(s) = -\beta \;-\; \lambda\,\mathrm{imp}(s) \;-\; \gamma\,[\mathrm{type}(s)\text{ is a catch-all}].
\end{equation}
$\beta$ is the cost of a segment. $\mathrm{imp}(s)$ counts the actions of $s$ outside
its most frequent macro-class (execution, edit, exploration, version control,
delegation). The last term prefers spans that the decision table
(Section~\ref{sec:types}) can type specifically.

\paragraph{Boundary score.}
\begin{equation}
b(c) = w_{p}\,\mathrm{PASS}(c) + w_{f}\,\mathrm{FAIL}(c)\,\mathrm{SW}(c) + w_{a}\,\mathrm{AREA}(c) + w_{s}\,[c{+}1 = A] - \infty\cdot\mathrm{HARD}(c).
\end{equation}
The terms are:
\begin{itemize}
  \item $\mathrm{PASS}(c)$: $a_c$ is a passing test run, so the goal was reached;
  \item $\mathrm{FAIL}(c)\,\mathrm{SW}(c)$: $a_c$ is a decisive failure \emph{and} the
        macro-class changes at $c$. A failure alone is not a boundary, because agents
        usually keep working on the same goal after a failure;
  \item $\mathrm{AREA}(c)$: two lookups on either side of $c$ concern different code
        areas;
  \item the fourth term rewards a cut just before the anchor.
\end{itemize}

\paragraph{Hard constraints.} $\mathrm{HARD}(c)$ forbids a cut that would separate:
\begin{itemize}
  \item an edit from its first decisive run (an attempt includes its judgement);
  \item a delegation from its sub-agent's actions;
  \item a poll from the run it polls;
  \item two runs of a strict repeat: the same command failing with the same error
        signature, with no edit or environment change between them.
\end{itemize}
Short interruptions (``blips'') are not a rule. They are absorbed whenever keeping them
costs less ($\lambda$ per off-class action) than the two extra segments that cutting
around them would create ($2\beta$).

\paragraph{Algorithm.} Equation~\ref{eq:objective} decomposes over spans and cuts, so it
is maximised exactly by the dynamic program of Algorithm~\ref{alg:dp}, the
optimal-partitioning recurrence~\cite{jackson2005optimal,killick2012pelt}. Because the
decision table fixes each span's type, semi-Markov segmentation with typed
spans~\cite{sarawagi2004semicrf} reduces to this one-dimensional form.
\begin{itemize}
  \item \textbf{Cost:} $O(n^2)$ spans, with $\mathrm{imp}$ in $O(1)$ from prefix counts
        and typing in $O(|s|)$, so $O(n^3)$ in the worst case. With $n\le117$ this is
        well under a second per trajectory.
  \item \textbf{Exactness:} the result is the global optimum of the stated objective.
        What remains heuristic is the objective's evidence and weights, which are
        explicit, attributable per boundary, and tested by ablation and sensitivity
        analysis (Section~\ref{sec:design}).
\end{itemize}

\begin{algorithm}[t]
\caption{Outcome-anchored optimal partitioning}
\label{alg:dp}
\begin{algorithmic}[1]
\Require actions $a_1..a_n$ with frames, results, anchor $A$; scores $\varphi$, $b$
\State $F(0) \gets 0$
\For{$j \gets 1$ \textbf{to} $n$}
  \State $F(j) \gets \max_{0\le i<j}\big[F(i) + \varphi(a_{i+1..j}) + [i>0]\,b(i)\big]$
  \State $\mathit{back}(j) \gets$ the maximising $i$
\EndFor
\State follow $\mathit{back}$ from $n$ to $0$ to recover the spans
\State \Return spans, each typed by the decision table (Section~\ref{sec:types})
\end{algorithmic}
\end{algorithm}

\paragraph{Example.} In Figure~\ref{fig:timeline}, actions 1--3 stay together: the
failed test run at action 1 is followed by a read of the harness script, which counts
as execution work, so there is no class switch. Actions 6--7 cannot be cut (an edit and
its judging run). A cut after action 7 is rewarded: a decisive failure followed by a
switch to exploration.

\subsection{Segment types and outcomes}
\label{sec:types}
A segment's type is decided by an ordered decision table. The \emph{type} encodes the
goal the segment served and what came of it; \emph{flags} encode how. The table asks,
in order:
\begin{enumerate}
  \item is it a stuck repetition?
  \item is it a delegation?
  \item does it edit application source? Then it is a fix attempt: passed, failed,
        breaks the build, blocked by the environment, or unverified;
  \item does a test execute? Then it is a reproduction (before the anchor) or a
        validation (after), succeeded or failed;
  \item is it a test or harness edit with no decisive run?
  \item is its main work making things run? Then it is blocked, or toolchain work;
  \item is it a review, an orientation, or a lookup (targeted, a file-by-file sweep,
        or failed tool calls)?
\end{enumerate}
The pre-/post-fix split in question 4 is where context enters the labels: a failing
test is \code{reproduction\_succeeded} before the anchor and \code{fix\_attempt\_failed}
or \code{validation\_failed} after it.

Making the type depend on the goal rather than on the files touched matters in
practice. The first version of the prototype, a cascade of rules with a precedence
list, typed five Gemini segments as toolchain work, although in each one the test
executed and failed before the fix. The decision table types them as reproductions
and records the toolchain edit as a flag.

\subsection{Purposes}
\label{sec:purpose}
The segments $s_1..s_K$ are labelled with seven purposes:
\code{orientation}, \code{localization}, \code{reproduction}, \code{environment},
\code{repair}, \code{validation} and \code{review}. We maximise
\begin{equation}
\sum_k e(s_k,p_k) - \delta\,\big|\{k: p_k\neq p_{k-1}\}\big|
\end{equation}
by Viterbi decoding~\cite{rabiner1989hmm}. The emission $e(s,p)$ is:
\begin{itemize}
  \item 0 when $p$ is the segment's own purpose;
  \item $-c(\mathrm{own}(s),p)\,|s|$ when the codebook allows $s$ to nest in $p$;
  \item $-\infty$ otherwise, and also for a reproduction after the anchor or a
        validation before it.
\end{itemize}
A segment of $m$ actions between two segments of the same purpose nests exactly when
$c\,m < 2\delta$. Three costs therefore reproduce the codebook's nesting rules:
\begin{itemize}
  \item short interruptions nest;
  \item lookups and re-runs inside a fix--verify loop belong to repair;
  \item environment work done to get the failing test running belongs to reproduction.
\end{itemize}

\subsection{Surfacing failure and inefficiency}
\label{sec:surface}
\paragraph{Flags.} Each segment carries flags:
\begin{itemize}
  \item \emph{process integrity}: edits to the test oracle, the harness or the
        toolchain;
  \item \emph{wasted effort}: a stuck repeated command, environment wrestling,
        re-reading the same file, exploration off the eventual patch, reverting
        changes, tool errors;
  \item \emph{what followed a failure}: retried immediately, retried after a detour,
        abandoned, or the trajectory ended unresolved.
\end{itemize}

\paragraph{Loops.} A loop overlay marks fix--verify loops, reproduction loops and
stuck segments.

\paragraph{Run badges.} Each run carries badges: whether the agent reproduced the bug
before fixing, edited the test or the harness, was blocked by the environment, and
whether its final message claims a pass that the harness contradicts.

\paragraph{Views.} Six coordinated views share one state, so a filter, a selection or a
playhead in one view updates the others:
\begin{enumerate}
  \item a collapsible top-down outline (Figure~\ref{fig:topdown});
  \item a timeline with purpose and segment bands, loop overlays and a failure rail
        (Figure~\ref{fig:timeline});
  \item a phase Sankey split by agent, by outcome, or by failure density;
  \item a category-by-step fingerprint for the three agents of one bug;
  \item an agent comparison aligned over segment types;
  \item statistics per agent.
\end{enumerate}
An ``only failure and inefficiency'' filter greys out everything else.

\begin{figure}[t]
  \centering
  \includegraphics[width=\columnwidth]{figs/traj-topdown.png}
  \caption{Top-down view of the same run: purposes (left bars) contain segments, each
  with type, outcome, flags and its action range. Segments that nest in another
  purpose are marked. Prototype output; actions 6--7 show the toolchain-typing error
  fixed by the decision table (Section~\ref{sec:types}).}
  \Description{A collapsible outline of purposes and their segments with outcome pills
  and flags.}
  \label{fig:topdown}
\end{figure}

%% ===================================================================== 4
\section{Experimental Design}
\label{sec:design}

\subsection{Research questions}
\begin{description}
  \item[RQ1, segmentation validity.] How accurately does \tool{} recover
        human-annotated segment boundaries, types, outcomes and purposes, compared with
        baselines from existing tools and with agreement between two humans?
  \item[RQ2, surfacing.] How precisely do its flags identify failure and inefficiency,
        compared with human judgement and with existing detectors?
  \item[RQ3, behavioural insight.] Which segment-level behaviours separate passing from
        failing runs, and how do the three agents differ?
  \item[RQ4, usefulness.] Do analysts find a run's failure cause faster and more
        accurately with \tool{} than with a step-list viewer?
\end{description}

\subsection{Corpus}
Fifty Android bugs, each attempted by Claude Code, Codex and Gemini CLI, in a setting
where the agent receives the bug report, one failing UI test and a few starting
locations (Table~\ref{tab:corpus}). Outcomes are the harness's own test results for
that setting. The three agents ran on different machines, so environment failures are
partly host effects; we analyse each agent's process separately and report host as a
threat.

\begin{table}[t]
\centering
\small
\caption{Corpus. Medians per trajectory; segments and purposes from the current
prototype.}
\label{tab:corpus}
\resizebox{\columnwidth}{!}{%
\begin{tabular}{@{}lrrrrrr@{}}
\toprule
\textbf{Agent} & \textbf{Runs} & \textbf{Actions} & \textbf{Med.\ act.} & \textbf{Med.\ seg.} & \textbf{Med.\ purp.} & \textbf{Pass} \\
\midrule
Claude Code & 50 & 1{,}322 & 18.5 & 5  & 3   & 40 \\
Codex       & 50 & 1{,}985 & 35.5 & 11 & 6   & 45 \\
Gemini CLI  & 50 & 2{,}093 & 35.5 & 13 & 6.5 & 39 \\
\midrule
Total       & 150 & 5{,}400 & 33 & 10 & 6 & 124 \\
\bottomrule
\end{tabular}}
\end{table}

\subsection{Reference set}
Two raters label segments, types, outcomes, flags and purposes with a codebook that
defines eight immediate goals, five keep rules and five cut rules in priority order, an
ordered type decision table, and three nesting rules for purposes, with a fully worked
example.
\begin{itemize}
  \item \textbf{Blind.} Raters see each action's tool, target, category, result and a
        result preview, never the tool's segments.
  \item \textbf{Pilot.} Bug 18 $\times$ 3 agents (99 actions) calibrates the raters
        and sets the segmenter's free weights.
  \item \textbf{Test set.} Five bugs $\times$ 3 agents (708 actions; 7 passing and 8
        failing runs), adjudicated into the reference.
  \item \textbf{Freeze.} Rules, weights and codebook are frozen before test labelling
        starts.
  \item \textbf{Audit.} The mid-level category and the result parser are deterministic
        and are not rated. One author audits a stratified sample of about 200 actions
        and reports precision per category and per result kind.
\end{itemize}

\subsection{Baselines}
\begin{itemize}
  \item \textbf{Phase runs}, after Graphectory~\cite{graphectory2026}: a phase tag per
        command, then run-length segments.
  \item \textbf{BPE procedures}, after ProcGrep~\cite{oderinwale2026procgrep}: $K\in\{10,20,40\}$
        merges, fitted leave-one-bug-out.
  \item \textbf{Trivial floors:} category run-length, fixed windows, and random
        boundaries with the same count.
  \item \textbf{Clustering:} $k$-means typing of segments.
  \item \textbf{A linear-chain CRF}~\cite{lafferty2001crf,tax2016eventabstraction}
        trained on the reference labels, leave-one-bug-out.
  \item \textbf{An LLM segmenter} that must cite step ids, alone and on top of our
        segments.
  \item \textbf{For RQ2,} re-implemented detectors: AgentDiff's stagnant
        loop~\cite{agentdiff}, Mindwalk's edits-after-last-verify~\cite{mindwalk},
        and a recovery rule from a Claude Code session visualiser.
\end{itemize}

\subsection{Metrics}
\begin{itemize}
  \item \textbf{Boundaries:} precision, recall and F1 with one-to-one matching at
        tolerance 0 and $\pm1$, and WindowDiff~\cite{pevzner2002windowdiff}.
  \item \textbf{Types and outcomes:} Cohen's $\kappa$~\cite{cohen1960kappa} on matched
        segments (IoU $\ge 0.5$).
  \item \textbf{Purposes:} per-action $\kappa$.
  \item \textbf{Flags:} precision and recall per flag.
  \item \textbf{Ablation of C1:} the outcome terms, then the state terms, removed.
  \item \textbf{Sensitivity:} each weight varied by $\pm25\%$ and $\pm50\%$.
  \item \textbf{Generality of C2:} leave-one-agent-out agreement of the frame
        classifier with tool names hidden, and coverage on trajectories of agents we
        never tuned for.
\end{itemize}

%% ===================================================================== 5
\section{Results}
\label{sec:results}

\subsection{Preliminary: what the segments reveal}
Before the reference set exists, the prototype's output already separates the agents
(Table~\ref{tab:behaviour}). These counts come from rules that are not yet validated;
RQ1 will establish how far they can be trusted.

\begin{table}[t]
\centering
\small
\caption{Runs (of 50 per agent) showing each behaviour, from the current prototype.}
\label{tab:behaviour}
\resizebox{\columnwidth}{!}{%
\begin{tabular}{@{}lrrr@{}}
\toprule
\textbf{Behaviour} & \textbf{Claude} & \textbf{Codex} & \textbf{Gemini} \\
\midrule
failing test run before the first source edit & 0  & 3  & 50 \\
$\ge1$ build/test run blocked by the environment & 2  & 50 & 18 \\
edits the test script                            & 0  & 1  & 14 \\
edits the harness script                         & 2  & 0  & 6  \\
final message claims tests pass                  & 45 & 44 & 35 \\
\quad \ldots{} while the harness reports failure & 8  & 0  & 5  \\
$\ge1$ fix--verify loop                          & 8  & 8  & 21 \\
$\ge1$ stuck segment                             & 3  & 3  & 1  \\
\bottomrule
\end{tabular}}
\end{table}

The agents differ in process, not just in pass rate:
\begin{itemize}
  \item \textbf{Gemini CLI} runs the failing test before editing source in every run,
        and edits the test script in 14.
  \item \textbf{Claude Code} never runs a test before its first source edit, and 8 of
        its 10 failing runs end with a claim that the tests pass.
  \item \textbf{Codex} meets an environment block in every run, a host effect that the
        flags make visible instead of hiding it inside the pass rate.
\end{itemize}

Existing detectors are nearly silent on the same corpus. Our re-implementations of
AgentDiff's stagnant-loop rule fire in 0\%, 6\% and 0\% of runs (Claude, Codex,
Gemini), and Mindwalk's edits-after-last-verify in 0\%, 0\% and 4\%.

\subsection{RQ1: segmentation validity}
\TODO{Boundary F1 (tol.\ 0/$\pm1$), WindowDiff, type/outcome $\kappa$ and purpose
$\kappa$ for \tool{}, every baseline, and the human ceiling; ablation without outcome
and without state evidence; sensitivity table.}

\subsection{RQ2: surfacing failure and inefficiency}
\TODO{Per-flag precision and recall against the reference set and ${\sim}30$ sampled
flag instances; comparison with the re-implemented detectors.}

\subsection{RQ3: behavioural insight}
\TODO{Segment-level behaviours vs.\ harness outcome: effect sizes, Fisher/Mann--Whitney
tests with Holm correction, same-bug contrasts. Exploratory: 26 failing runs.}

\subsection{RQ4: usefulness}
\TODO{Analyst study, 5--8 participants, failure-cause tasks: accuracy, time, SUS.}

%% ===================================================================== 6
\section{Discussion and Threats to Validity}
\label{sec:discussion}

\paragraph{Rules versus learning.} \tool{}'s evidence and weights are designed, not
learned. With 708 labelled actions, a learned segmenter (our CRF baseline) has little to
learn from. A stated objective with about ten weights can be inspected, ablated and
fitted on a pilot. Exact decoding removes the order effects that made the first
prototype mislabel segments.

\paragraph{Shared specification.} Raters and tool share the codebook's definitions,
including its default sizes, by design. RQ1 therefore measures how well the tool
applies the rubric, while agreement between the two raters measures how ambiguous the
rubric is.

\paragraph{Internal validity.}
\begin{itemize}
  \item The pilot that sets the weights is small (99 actions); the sensitivity analysis
        reports how fragile each weight is.
  \item The author who built the tool is also a rater; the freeze precedes test
        labelling.
\end{itemize}

\paragraph{External validity.}
\begin{itemize}
  \item One benchmark (Android repair), one task setting, and one run per bug and agent.
  \item Hosts differ per agent.
  \item Result parsing needs an adapter per harness.
  \item Gemini's delegated sub-agent work is not logged.
\end{itemize}
The leave-one-agent-out and unseen-agent checks address generality of the input layer
only.

%% ===================================================================== 7
\section{Related Work}
\label{sec:related}

\paragraph{Trajectory viewers and observability.}
\begin{itemize}
  \item \emph{Step views:} Agent Trajectory Explorer renders thought--action--observation
        turns~\cite{desmond2025explorer}; SeaView, TraceView and AgentStepper support
        experiment browsing, manual labelling and stepping through
        runs~\cite{seaview2025,traceview2026,agentstepper2026}.
  \item \emph{Per-step phases and mining:} Graphectory tags actions with phases and
        mines phase strings of resolved and unresolved runs~\cite{graphectory2026};
        Agent ATO tags interactions with four rule-based
        kinds~\cite{agentato2026}.
  \item \emph{Corpus-scale analysis:} visual analytics for tree-search coding
        agents~\cite{wang2025illuminating}, Docent~\cite{docent2025} and Inspect
        Scout~\cite{inspectscout} summarise or scan transcripts with LLMs.
\end{itemize}
\tool{} adds contiguous, outcome-anchored segments between the step and the run, and
validates them.

\paragraph{Studies of agent trajectories and failures.}
\begin{itemize}
  \item \emph{Trajectory studies:} Bouzenia and Pradel annotate thought--action--result
        trajectories of repair agents~\cite{bouzenia2025tar}; recent studies analyse
        failure as a temporal process~\cite{zhao2026failureprocess}, ``lucky''
        passes~\cite{sahoo2026agentlens} and cost-wasting
        behaviours~\cite{hu2026costinefficient}.
  \item \emph{Failure taxonomies and attribution:} MAST and TRAIL contribute taxonomies
        and annotated traces~\cite{cemri2025mast,trail2025}; CHIEF, FALAT, TrajAudit
        and HarnessFix attribute failures to steps, agents or harness
        components~\cite{chief2026,rafi2026falat,trajaudit2026,chen2026harnessfix}.
\end{itemize}
Our segments are a natural unit for such attribution.

\paragraph{Structure inside trajectories.}
\begin{itemize}
  \item \emph{Fixed workflows and stages:} Agentless fixes a localize--repair--validate
        workflow~\cite{xia2025agentless}, and Kimi-Dev annotates agent trajectories
        with stages using an LLM~\cite{kimidev2025}.
  \item \emph{Induced units:} ProcGrep induces procedures by
        BPE~\cite{oderinwale2026procgrep}, behavioural ``genomes'' encode actions over a
        small alphabet~\cite{deng2026genome}, and Agent Workflow Memory induces
        reusable routines~\cite{wang2025awm}.
  \item \emph{Agent-declared boundaries:} Wei has agents declare segment boundaries
        while they work~\cite{wei2026selfsegment}, which cannot be applied to existing
        logs.
\end{itemize}

\paragraph{Segmentation and event abstraction.}
\begin{itemize}
  \item \emph{Text segmentation:} cohesion and optimal
        partitioning~\cite{hearst1997texttiling,utiyama2001segmentation}, evaluated
        with WindowDiff~\cite{pevzner2002windowdiff}.
  \item \emph{Change-point detection:} exact dynamic programs for optimal
        partitions~\cite{jackson2005optimal,killick2012pelt}.
  \item \emph{Structured prediction:} semi-Markov models that label
        spans~\cite{sarawagi2004semicrf,lafferty2001crf}.
  \item \emph{Event abstraction in process mining:} mapping low-level events to
        activities~\cite{vanzelst2021eventabstraction,mannhardt2018guided,tax2016eventabstraction}.
\end{itemize}
\tool{} is an instance of event abstraction for agent logs. What is specific to it is
the evidence: run results and the pre-/post-fix state.

%% ===================================================================== 8
\section{Conclusion}
\label{sec:conclusion}
\tool{} segments coding-agent trajectories into contiguous units with goals, types and
outcomes, using an exact optimisation over evidence that existing tools ignore: what
each run reported, and whether the fix had begun. On 150 trajectories of three agents,
the segments already separate the agents' repair processes and expose failures that
existing detectors miss. \TODO{one sentence of RQ1--RQ2 results.} We release the
converters, the segmenter, the viewer, the codebook and the reference set.

\balance
\bibliographystyle{ACM-Reference-Format}
\bibliography{paper}

\end{document}
